% !TeX root = ../main.tex
\chapter{Continuità}
\label{cap:continuita}
\fonti{Lezioni 13 e 14}
% Fonte: Lezione 13, pp. 13--17. Prova dell'algebra e chiarimenti su estremi/isolati aggiunti.

\section{Continuità e successioni}

Una funzione è continua in un punto quando avvicinare gli argomenti a quel punto avvicina anche i valori della funzione al valore assunto nel punto. Le successioni permettono di esprimere questa idea usando la nozione di limite già introdotta.

\begin{definizione}[Continuità in un punto e su un insieme]
\label{def:continuita-sequenziale}
Siano $f\colon A\subseteq\R\to\R$ e $x_0\in A$. Si dice che $f$ è \emph{continua in $x_0$} se, per ogni successione $(x_n)$ a valori in $A$,
\[
x_n\longrightarrow x_0
\quad\Longrightarrow\quad
f(x_n)\longrightarrow f(x_0).
\]
Se è continua in ogni punto di $A$, si dice \emph{continua su $A$} e si scrive $f\in C(A)$ oppure $f\in C^0(A)$.
\end{definizione}

La condizione riguarda tutte le successioni del dominio che tendono a $x_0$, comprese quelle che assumono il valore $x_0$. Il punto deve appartenere al dominio, perché compare $f(x_0)$; la notazione $C^0$ non richiede l'esistenza di derivate.

\begin{esempio}[Il valore assoluto]
Per ogni $x_0\in\R$, la proposizione~\ref{prop:limite-modulo} dà $x_n\to x_0\Rightarrow|x_n|\to|x_0|$. Pertanto $x\mapsto|x|$ è continua su tutta $\R$, compreso $0$.
\end{esempio}

\begin{schema}
\begin{tikzpicture}[x=1.6cm,y=1cm]
  \draw[asse] (-.3,0) -- (5,0) node[right] {$x$};
  \draw[asse] (0,-.2) -- (0,3.4) node[above] {$y$};
  \draw[intervallo,domain=.45:4.5,samples=80] plot (\x,{1+.35*\x+.06*(\x-2)^3});
  \draw[dashed,black!55] (3.5,0) -- (3.5,2.4275) -- (0,2.4275);
  \fill[coloreschema] (3.5,2.4275) circle[radius=2pt];
  \foreach \a in {2.8,3.1,3.3,3.4} {
    \pgfmathsetmacro{\b}{1+.35*\a+.06*(\a-2)^3}
    \draw[black!40] (\a,0) -- (\a,\b) -- (0,\b);
    \fill[coloreschema] (\a,\b) circle[radius=1.6pt];
  }
  \node[below] at (3.5,0) {$x_0$};
  \node[below] at (2.8,-.15) {$x_n$};
  \draw[asse,coloreschema] (2.85,-.55) -- (3.4,-.55);
  \node[left] at (0,2.4275) {$f(x_0)$};
  \node[left] at (0,1.7) {$f(x_n)$};
\end{tikzpicture}
\par\smallskip Schema di continuità: ai punti $x_n$ che si avvicinano a $x_0$ corrispondono ordinate che si avvicinano a $f(x_0)$. Il disegno mostra una sola scelta di successione; la definizione le richiede tutte.
\end{schema}

\section{Il ruolo del dominio}
% Fonte: Lezione 13, p. 16; preservati dominio [0,3] unito a {4}, valore 5 e x_n=3-1/n.

La continuità si valuta usando soltanto punti del dominio. A un estremo di un intervallo le successioni possono avvicinarsi soltanto dal lato disponibile; un punto isolato presenta una situazione ancora diversa.

\begin{esempio}[Un estremo e un punto isolato]
Consideriamo una funzione $f\colon[0,3]\cup\{4\}\to\R$, continua sul tratto $[0,3]$, e poniamo $f(4)=5$. La successione $x_n=3-1/n$, per $n\ge1$, appartiene al dominio e tende a $3$. La continuità in $3$ richiede $f(3-1/n)\to f(3)$; non impone che questi valori tendano a $f(4)=5$, perché $x_n$ non tende a $4$.

Il punto $4$ è isolato nel dominio. Se $x_n\in[0,3]\cup\{4\}$ e $x_n\to4$, usando il raggio $1/2$ si ha definitivamente $x_n\in(7/2,9/2)$. L'unico punto del dominio in tale intervallo è $4$, quindi $x_n=4$ e $f(x_n)=5$ definitivamente. La funzione è continua anche in $4$, qualunque sia il valore che le si assegna in quel punto.
\end{esempio}

\begin{schema}
\begin{tikzpicture}[x=1.65cm,y=.7cm]
  \draw[asse] (-.3,0) -- (5,0) node[right] {$x$};
  \draw[asse] (0,-.2) -- (0,5.8) node[above] {$y$};
  \draw[intervallo,domain=0:3,samples=70] plot (\x,{1+\x/3+\x*(\x-1.5)*(\x-3)/5});
  \fill[coloreschema] (0,1) circle[radius=2pt];
  \fill[coloreschema] (3,2) circle[radius=2pt];
  \fill[coloreschema] (4,5) circle[radius=2pt];
  \draw[dashed,black!50] (4,0) -- (4,5) -- (0,5);
  \draw[dotted,black!50] (3,0) -- (3,2);
  \foreach \a in {0,3,4} {\node[below] at (\a,0) {$\a$};}
  \node[left] at (0,5) {$5$};
  \node[right] at (3,2) {$f(3)$};
  \draw[asse,coloreschema] (2.05,.12) -- (2.95,.12);
  \node[below] at (2.3,-.25) {$3-1/n$};
\end{tikzpicture}
\par\smallskip Fra $3$ e $4$ il dominio non contiene punti: il grafico non deve essere congiunto attraverso questo intervallo. La continuità nel punto isolato $4$ non dipende dal valore all'estremo $3$.
\end{schema}

Il ragionamento sul punto $4$ vale in generale: se $x_0$ è isolato in $A$, ogni successione di $A$ che tende a $x_0$ è definitivamente uguale a $x_0$. Ogni funzione $A\to\R$ è dunque continua nei punti isolati del proprio dominio.

\section{Funzioni elementari e operazioni}
% Fonte: Lezione 13, p. 17. Primo teorema enunciato senza prova nella fonte; prova del secondo aggiunta.

\begin{teorema}[Continuità delle funzioni elementari]
\label{teo:continuita-elementari}
Le funzioni elementari già introdotte, potenze, funzioni trigonometriche, esponenziali, logaritmi e funzioni iperboliche, e le loro inverse sui rami in cui sono definite, sono continue nei rispettivi domini.\footnote{Questo risultato è assunto qui come proprietà delle funzioni elementari. Non è una conseguenza della sola definizione: le dimostrazioni generali, in particolare per trigonometria ed esponenziali, non sono sviluppate in questa trattazione. Restano le ipotesi dichiarate nella costruzione dell'esponenziale reale.}
\end{teorema}

L'affermazione riguarda soltanto punti del dominio: per esempio, il logaritmo è continuo per argomenti positivi; la tangente nei punti in cui il coseno non si annulla. Per le radici di indice pari la continuità in $0$ si valuta nel dominio $[0,+\infty)$.

\begin{teorema}[Algebra delle funzioni continue]
Siano $f,g\colon A\subseteq\R\to\R$ continue in $x_0\in A$. \begin{enumerate}[label=\arabic*.]
\item Allora $f+g$, $f-g$ e $fg$ sono continue in $x_0$.
\item Se inoltre $g(x_0)\ne0$, anche $f/g$, definita su $B=\{x\in A:g(x)\ne0\}$, è continua in $x_0$.
\end{enumerate}
\end{teorema}
\begin{proof}
\begin{enumerate}[label=\arabic*.]
\item \textbf{Somma, differenza e prodotto.} Per ogni successione $x_n\in A$ tendente a $x_0$, la continuità dà $f(x_n)\to f(x_0)$ e $g(x_n)\to g(x_0)$. L'algebra dei limiti finiti implica
\[
\begin{aligned}
f(x_n)\pm g(x_n)&\longrightarrow f(x_0)\pm g(x_0),\\
f(x_n)g(x_n)&\longrightarrow f(x_0)g(x_0).
\end{aligned}
\]
Sono esattamente le condizioni di continuità delle tre funzioni in $x_0$.
\item \textbf{Quoziente.} Poiché $g(x_0)\ne0$, si ha $x_0\in B$. Se $x_n\in B$ e $x_n\to x_0$, valgono gli stessi limiti dei valori di $f$ e $g$; il teorema~\ref{teo:limite-quoziente} dà
\[
\frac{f(x_n)}{g(x_n)}\longrightarrow\frac{f(x_0)}{g(x_0)}.
\]
Questa è la continuità di $f/g$ in $x_0$ sul suo dominio $B$.
\end{enumerate}
\end{proof}

\begin{esempio}
Ogni polinomio è continuo su $\R$, perché è una somma finita di costanti e potenze. Ogni funzione razionale $P/Q$ è continua dove $Q$ non si annulla. Per esempio $1/(x^2-4)$ è continua su $\R\setminus\{-2,2\}$: il risultato non assegna valori nei due punti esclusi.
\end{esempio}

% !TeX root = ../main.tex
% Fonte: Lezione 14, pp. 2--6. Accorpati i richiami alla definizione; prova della composizione aggiunta.
\section{Composizione e calcolo dei limiti}

\begin{proposizione}[Continuità della composta]
Siano $f\colon A\to B$ e $g\colon B\to\R$, con $A,B\subseteq\R$. Se $f$ è continua in $x_0\in A$ e $g$ è continua in $f(x_0)$, allora $g\circ f$ è continua in $x_0$.
\end{proposizione}
\begin{proof}
Per ogni successione $x_n\in A$ con $x_n\to x_0$, la continuità di $f$ dà $f(x_n)\to f(x_0)$. I valori $f(x_n)$ appartengono a $B$; la continuità di $g$ nel punto $f(x_0)$ implica quindi $g(f(x_n))\to g(f(x_0))$, come richiesto.
\end{proof}

La stessa definizione si applica direttamente al calcolo di un limite: se $x_n\to x_0$, i termini appartengono al dominio di $f$ e $f$ è continua in $x_0$, allora $f(x_n)\to f(x_0)$. Si può dunque calcolare prima il limite dell'argomento, controllare dominio e continuità, e poi applicare la funzione al limite trovato. Per una successione definita solo definitivamente si lavora sulla sua coda.

\begin{esercizio}[Continuità e fattore oscillante]
Calcolare, per $n\to+\infty$,
\[
\sin\!\left(\left|2\left[\cos\!\left(\frac1{n^2}\right)-1\right]\sin n-\pi\right|\right).
\]
\end{esercizio}
\begin{svolgimento}
Poiché $1/n^2\to0$ e il coseno è continuo, $\cos(1/n^2)\to1$. Quindi $2[\cos(1/n^2)-1]\to0$. La successione $\sin n$ è limitata, anche se non ha limite, e il prodotto di una limitata per un'infinitesima tende a zero:
\[
2[\cos(1/n^2)-1]\sin n-\pi\longrightarrow-\pi.
\]
Per la continuità del modulo, i valori assoluti tendono a $\pi$. Infine, per la continuità del seno, la successione richiesta tende a $\sin\pi=0$.
\end{svolgimento}

\begin{esercizio}[Una composizione di funzioni elementari]
Calcolare
\[
\lim_{n\to+\infty}\sqrt[3]{\cos(1/n+\sin(1/n^3))}.
\]
\end{esercizio}
\begin{svolgimento}
Si procede dall'argomento più interno verso l'esterno. Poiché $1/n^3\to0$ e il seno è continuo, $\sin(1/n^3)\to0$. Sommando $1/n\to0$ si ottiene $1/n+\sin(1/n^3)\to0$. La continuità del coseno dà
\[
\cos(1/n+\sin(1/n^3))\longrightarrow\cos0=1.
\]
La radice cubica è definita e continua su tutta $\R$, perciò il limite cercato è $\sqrt[3]1=1$.
\end{svolgimento}

\section{Quando la sostituzione non è lecita}

\begin{esempio}[Una funzione discontinua]
Consideriamo
\[
f(x)=\begin{cases}1,&x=0,\\0,&x\ne0.\end{cases}
\]
Per $n\ge1$, $f(1/n)=0$, quindi $f(1/n)\to0$. Sebbene $1/n\to0$, il limite non è $f(0)=1$: manca la continuità in $0$. La successione costante $x_n=0$ dà invece $f(x_n)=1$, confermando che le immagini di successioni tendenti allo stesso punto possono avere limiti diversi.
\end{esempio}

\begin{schema}
\begin{tikzpicture}[x=1.6cm,y=1.5cm]
  \draw[asse] (-2.2,0) -- (2.4,0) node[right] {$x$};
  \draw[asse] (0,-.3) -- (0,1.55) node[above] {$y$};
  \draw[intervallo] (-2,0) -- (2.1,0);
  \draw[estremo] (0,0) circle[radius=2pt];
  \fill[coloreschema] (0,1) circle[radius=2pt] node[left] {$1$};
  \node[below left] at (0,0) {$0$};
\end{tikzpicture}
\par\smallskip Il punto pieno indica $f(0)=1$; il punto vuoto esclude $(0,0)$ dal grafico. Avvicinarsi a zero attraverso punti non nulli dà sempre valore $0$.
\end{schema}

\begin{esempio}[Controllare prima il dominio]
La formula $\sqrt{\cos(1/n^2)-1}$, per $n\ge1$, non definisce una successione reale: infatti $0<1/n^2\le1<2\pi$, dunque $\cos(1/n^2)<1$ e il radicando è negativo. Il fatto che il radicando tenda a zero non rende definita la radice quadrata dei suoi termini. Non è quindi lecito concludere che la successione abbia limite $\sqrt0$ in $\R$.
\end{esempio}

